\documentclass{article}
\usepackage[T1]{fontenc}
\usepackage[utf8]{inputenc}
\usepackage{graphicx}
\usepackage{wrapfig}
\usepackage{multicol}
\usepackage{booktabs}
\begin{document}

\title{Complex Layout with Image}
\author{Fixture Author}
\date{\today}
\maketitle

\section{Introduction to the Complex Layout}
This fixture is the Stage 3 example: a genuinely complex document layout that
combines a title block, section headings, an embedded raster image, text
wrapping around that image, a multi-column region, and a table all in one
document. Its primary assertion target is the section heading
\textbf{Complex Layout with Image} and the prose marker
\texttt{layoutsentinel-intro}, which appears only here and nowhere else in the
document.

\begin{wrapfigure}{r}{0.4\textwidth}
  \centering
  \includegraphics[width=0.35\textwidth]{sample.png}
  \caption{Sentinel image caption for the wrapped figure.}
  \label{fig:sample}
\end{wrapfigure}

The image above is embedded as a solid-color square with no text content. The
surrounding prose, however, is meant to be extracted, and it carries the
distinctive marker \texttt{layoutsentinel-wrap}. The \texttt{wrapfig}
environment places the figure on the right and lets this paragraph flow around
it, so the physical layout here is more complex than the plain single-column
fixtures. When this document is parsed, the text of this paragraph should still
be recovered even though the figure sits beside it.

Additional filler prose follows so that the wrapping behavior is visually
obvious and the paragraph genuinely extends beside the figure rather than
merely sitting above or below it. This describes the seasonal variation of
precipitation across the northern uplands, the gradual retreat of the glaciers,
and the corresponding shift in the distribution of alpine vegetation. None of
this filler carries a sentinel; it exists purely to give the layout real body
so the test fixture reads as an authentic complex page rather than a toy.

\section{Multi-Column Region}
\begin{multicols}{2}
In this section the text flows across two columns. The left-hand content
begins here and continues into the second column automatically. The marker
\texttt{layoutsentinel-multicol} is placed in this two-column region so that
the tests can confirm that multi-column text is still extracted as readable
prose, just as in the earlier two-column fixture. Additional narrative about
the measurement campaign, the calibration of the instruments, and the number of
recordings taken per season follows to give the columns reasonable height.

The second column carries the marker \texttt{layoutsentinel-multicol-2}. It
repeats the theme of the previous column while adding further detail about the
data processing pipeline: the raw readings are smoothed, outliers are flagged,
and the resulting series are aggregated into the summary table shown below.
Together these two sentinels let the tests verify that text spanning a
multi-column break is recovered intact.
\end{multicols}

\section{Summary Table}
The document also embeds a \texttt{tabular} table so that table extraction can
be exercised alongside the image and multi-column features. The table below
carries the sentinel \texttt{layoutsentinel-table} in its first data cell.

\begin{center}
\begin{tabular}{lcc}
\toprule
Region & Recordings & Flag \\
\midrule
layoutsentinel-table & 128 & active \\
Northern & 64 & active \\
Alpine & 32 & retired \\
\bottomrule
\end{tabular}
\end{center}

\section{Conclusion}
This concluding section spans the second page of the fixture and re-states the
heading \textbf{Complex Layout with Image} alongside the marker
\texttt{layoutsentinel-conclusion}. It ensures the document runs to at least
two pages and that a reliable, text-only assertion target exists on the final
page regardless of how the image is routed during extraction. The image itself
is blank, so the automated tests deliberately assert only on the surrounding
text and heading.

\end{document}
